\documentclass[10pt]{beamer}
\usetheme{metropolis} % Nowoczesny, czysty motyw
\usepackage[utf8]{inputenc}
\usepackage[T1]{fontenc}
\usepackage[polish]{polski}

\title{Projekt ZPI 2025: Innowacyjna Gra Roguelike}
\subtitle{Mechanika parowania jako klucz do sukcesu}
\date{10 kwietnia 2026}
\author{Kacper Trzciński, Aleksander Bielicki, Radosław Głasek, \\ Kuba Pączka, Michał Popkowicz, Filip Shramko}
\institute{Zarządzanie Przedsięwzięciami Informatycznymi}

\begin{document}

\maketitle

% Slajd 2: Problem i Rozwiązanie
\begin{frame}{Problem i Rozwiązanie}
    \textbf{Problem:} Rynek gier roguelike jest nasycony tytułami opartymi na czystej agresji i statystykach. Gracze szukają nowych wyzwań intelektualnych w dynamicznej oprawie. \par
    \vspace{0.5cm}
    \textbf{Rozwiązanie:} Nasza gra przesuwa ciężar z ataku na obronę. Centralnym punktem jest \textbf{mechanika parowania}, która służy nie tylko do ochrony, ale jest głównym narzędziem rozwiązywania zagadek i pokonywania przeciwników[cite: 9, 37].
\end{frame}

% Slajd 3: Unique Value Proposition (UVP)
\begin{frame}{Unique Value Proposition (UVP)}
    Co nas wyróżnia?
    \begin{itemize}
        \item \textbf{Defensywa to Ofensywa:} Gracz nie posiada klasycznej broni. Jego siła zależy od umiejętności przejęcia momentum przeciwnika[cite: 10, 37].
        \item \textbf{Logika w walce:} Połączenie gry zręcznościowej z zagadkami logicznymi (np. odbijanie pocisków w celu aktywacji mechanizmów)[cite: 67, 69].
        \item \textbf{Metodyka PRINCE2:} Gwarancja dostarczenia produktu dzięki rygorystycznemu podejściu do zarządzania jakością i zakresem[cite: 1, 15].
    \end{itemize}
\end{frame}

% Slajd 4: Analiza Interesariuszy
\begin{frame}{Analiza Interesariuszy}
    Projekt angażuje kluczowe role zgodnie z PRINCE2[cite: 5]:
    \begin{itemize}
        \item \textbf{Przewodnicząca Komitetu Sterującego:} Aleksandra Sobótka – nadzór strategiczny i biznesowy[cite: 5].
        \item \textbf{Główny Użytkownik:} Aleksandra Sobótka – definiowanie wymagań jakościowych[cite: 5].
        \item \textbf{Główny Dostawca:} Aleksander Płocharski – odpowiedzialny za jakość techniczną[cite: 5].
        \item \textbf{Zewnętrzni:} Gracze szukający gier indie oraz społeczność akademicka.
    \end{itemize}
\end{frame}

% Slajd 5: Harmonogram i Kamienie Milowe
\begin{frame}{Harmonogram - Droga do celu}
    Realizacja podzielona na kluczowe etapy[cite: 75]:
    \begin{itemize}
        \item \textbf{Kwiecień 2026:} Zamknięcie dokumentacji (Karta projektu, WBS, Harmonogram).
        \item \textbf{Październik 2026:} Sprint 1 – Prototyp mechaniki parowania i fizyki momentum.
        \item \textbf{Grudzień 2026:} Sprint 3 – Wdrożenie zróżnicowanych przeciwników i zagadek.
        \item \textbf{Styczeń 2027:} Sprint 5 – Finalizacja wersji grywalnej i testy.
        \item \textbf{Semestr Letni 2027:} Prezentacja przed nowym rocznikiem[cite: 12, 70].
    \end{itemize}
\end{frame}

% Slajd 6: Budżet i Zasoby
\begin{frame}{Budżet: Efektywność kosztowa}
    Nasz budżet wynosi \textbf{0 PLN} dzięki optymalizacji zasobów[cite: 51]:
    \begin{itemize}
        \item \textbf{Narzędzia:} Godot/Unity/Unreal Engine, GitHub, ProjectLibre – darmowe licencje [cite: 46-50].
        \item \textbf{Zasoby ludzkie:} 6-osobowy zespół ekspertów (programiści, grafik, projektant poziomów, tester) [cite: 49, 53-57, 80-84].
        \item \textbf{Kapitał czasu:} 180 roboczogodzin zainwestowanych w rozwój innowacji[cite: 49].
    \end{itemize}
\end{frame}

% Slajd 7: Zarządzanie Ryzykiem
\begin{frame}{Zarządzanie Ryzykiem}
    Zidentyfikowaliśmy i zaplanowaliśmy reakcje na kluczowe zagrożenia[cite: 75]:
    \begin{itemize}
        \item \textbf{Zbalansowanie mechaniki:} Wczesne prototypowanie w Sprincie 1 i iteracyjne testy.
        \item \textbf{Obciążenie zespołu:} Ścisła priorytetyzacja zadań i gotowość do redukcji zakresu (MoSCoW).
        \item \textbf{Brak pomysłów na design:} Wczesny brainstorming i stworzenie katalogu interakcji (momentum, odbicia)[cite: 79].
    \end{itemize}
\end{frame}

% Slajd 8: Dlaczego warto zainwestować?
\begin{frame}{Co zyska Inwestor?}
    Inwestycja w nasz projekt to:
    \begin{itemize}
        \item \textbf{Gotowy model biznesowy:} Sprawdzone ramy operacyjne (BMC) i jasna struktura ról.
        \item \textbf{Minimalne ryzyko finansowe:} Wykorzystanie darmowej infrastruktury i ogromny wkład własny zespołu.
        \item \textbf{Potencjał skalowania:} Możliwość rozbudowy o tryb wieloosobowy lub porty mobilne w przyszłości (obecnie poza zakresem)[cite: 40, 44].
    \end{itemize}
    \vspace{0.3cm}
    \centering \large \textbf{Budujemy grę, w którą sami chcemy grać!}
\end{frame}

\end{document}